% !TeX program = xelatex
\documentclass[12pt,a4paper]{article}

% ==================== 中文支持 ====================
\usepackage[UTF8]{ctex}

% ==================== 页面布局 ====================
\usepackage{geometry}
\geometry{left=2.5cm,right=2.5cm,top=2.5cm,bottom=2.5cm}

% ==================== 字体与颜色（Linux 通用版） ====================
\usepackage{xcolor}
% \usepackage{fontspec} % ctex 已经加载了 fontspec，无需重复
% 英文字体：Liberation Serif（开源，与 Times New Roman 兼容）
% \setmainfont{Liberation Serif} % 不指定，使用默认罗马字体
% 中文字体：Noto Serif CJK SC（Google 开源字体，多数 Linux 发行版已预装）
% \setCJKmainfont{Noto Serif CJK SC} % 不指定，让 ctex 自动选择
% 如果上述字体不存在，可改为其他可用字体，例如：
% \setCJKmainfont{WenQuanYi Micro Hei}

% ==================== 常用宏包 ====================
\usepackage{amsmath, amssymb}
\usepackage{graphicx}
\usepackage{float}
\usepackage{booktabs}
\usepackage{array}
\usepackage{enumitem}
\usepackage{hyperref}
\hypersetup{
    colorlinks=true,
    linkcolor=blue,
    urlcolor=blue,
    citecolor=blue
}
\usepackage{caption}
\usepackage{subcaption}
\usepackage{listings}
\usepackage{verbatim}
\usepackage{fancyhdr}
\usepackage{titlesec}
\usepackage{parskip}

% ==================== 代码块样式 ====================
\lstset{
    basicstyle=\ttfamily\small,
    keywordstyle=\color{blue},
    commentstyle=\color{gray},
    stringstyle=\color{red},
    numbers=left,
    numberstyle=\tiny\color{gray},
    stepnumber=1,
    breaklines=true,
    frame=single,
    tabsize=4,
    showstringspaces=false,
    captionpos=b,
    xleftmargin=2em,
    xrightmargin=2em
}

% ==================== 段落格式 ====================
\setlength{\parindent}{0pt}
\setlength{\parskip}{0.5\baselineskip}

% ==================== 页眉页脚 ====================
\pagestyle{fancy}
\fancyhf{}
\fancyhead[L]{\leftmark}
\fancyhead[R]{\thepage}
\renewcommand{\headrulewidth}{0.4pt}

% ==================== 标题格式 ====================
\titleformat{\section}{\Large\bfseries}{\thesection}{1em}{}
\titleformat{\subsection}{\large\bfseries}{\thesubsection}{1em}{}

% ==================== 文档信息 ====================
\title{\textbf{通用文章模板（Linux 适配版）}}
\author{作者姓名}
\date{\today}

% ==================== 正文开始 ====================
\begin{document}

\maketitle

\begin{abstract}
这里是文章的摘要。
\end{abstract}

\tableofcontents
\newpage

\section{引言}
这是引言内容。

\subsection{子节示例}
内容示例。

\section{核心内容}
\label{sec:theory}
\subsection{列表环境}
\begin{itemize}
    \item 第一项
    \item 第二项
\end{itemize}

\begin{enumerate}
    \item 第一步
    \item 第二步
\end{enumerate}

\subsection{表格}
\begin{table}[H]
    \centering
    \caption{示例表格}
    \begin{tabular}{lcc}
        \toprule
        名称 & 数值1 & 数值2 \\
        \midrule
        项目 A & 10 & 20 \\
        项目 B & 30 & 40 \\
        \bottomrule
    \end{tabular}
    \label{tab:example}
\end{table}

\subsection{数学公式}
行内公式：$E = mc^2$。

独立公式：
\begin{equation}
    \int_{0}^{\infty} e^{-x^2} dx = \frac{\sqrt{\pi}}{2}
    \label{eq:gaussian}
\end{equation}

\subsection{插入图片}
\begin{figure}[H]
    \centering
    \includegraphics[width=0.6\textwidth]{example-image}
    \caption{示例图片}
    \label{fig:example}
\end{figure}

\subsection{代码块}
\lstset{language=Python}
\begin{lstlisting}[caption={Python 示例代码}, label=code:python]
def hello():
    print("Hello, LaTeX!")
\end{lstlisting}

\subsection{引用}
如图 \ref{fig:example} 所示，表 \ref{tab:example} 列出了数据，公式 (\ref{eq:gaussian}) 是一个著名积分。

\section{结论}
总结全文。

\begin{thebibliography}{99}
\bibitem{ref1} 作者, \textit{书名}, 出版社, 年份.
\end{thebibliography}

\appendix
\section{附录标题}
附录内容。

\end{document}
